\documentclass[11pt]{article}
\usepackage[margin=1in]{geometry}
\usepackage{amsmath,amssymb,amsthm,booktabs,hyperref}
\newtheorem{theorem}{Theorem}
\newtheorem{lemma}{Lemma}
\title{A finite obstruction to integrality of Schmidt-type coefficients}
\author{Principia Math In Progress\\\small AI-assisted research note; author review only}
\date{27 September 2026}
\begin{document}
\maketitle
\begin{abstract}
We disprove Guo and Zeng's finite integrality conjecture by taking
$s=2$ and $N=14$. Three explicit modular left-kernel vectors give
incompatible necessary conditions on every exponent $r\ge6$.
The three smaller permitted exponents are excluded by exact fractions.
All finite certificates are specified here and checked by a short
standard-library Python program supplied with this note.
\end{abstract}

\section*{Task 1: the source question and its resolution}
Write
\[
 B_{n,k}=\binom nk\binom{n+k}{k},\qquad
 F_n(r)=\sum_{k=0}^n B_{n,k}^{\,r}.
\]
For positive integers $r,s$, the rational coefficients $c_k^{(r,s)}$
are uniquely determined by
\begin{equation}\label{eq:def}
 F_n(r)=\sum_{k=0}^n B_{n,k}^{\,s}c_k^{(r,s)}\quad(n\ge0).
\end{equation}
Uniqueness follows from the nonzero diagonal $B_{n,n}^{\,s}$.
Conjecture 4.1 of the final published paper~\cite{GZ}, also recorded as
\href{https://mathdb.com/p/323296}{MathDB 323296}, asks whether for every
$s>1$ and $N\ge0$ some integer $r>s$ makes all $c_0^{(r,s)},\ldots,c_N^{(r,s)}$
integers.

\begin{theorem}\label{thm:main}
For every integer $r>2$, at least one of
$c_0^{(r,2)},\ldots,c_{14}^{(r,2)}$ is not an integer.
Consequently Conjecture 4.1 is false.
\end{theorem}

\subsection*{Three necessary congruences}
For $d\ge0$ put $A_d=(B_{n,k}^2)_{0\le n,k\le d}$, interpreting
$B_{n,k}=0$ when $k>n$. If an integer row vector $v$ satisfies
$vA_d\equiv0\pmod m$, then integrality of $c_0^{(r,2)},\ldots,c_d^{(r,2)}$
implies
\begin{equation}\label{eq:necessary}
 \sum_{n=0}^d v_n F_n(r)\equiv0\pmod m.
\end{equation}
Indeed this is obtained by multiplying the first $d+1$ identities
in~\eqref{eq:def} by $v$. Consider the following vectors, with coordinates
indexed from zero:
\begin{align*}
v={}&(54,0,189,84,0,564,414,243,153,333,486,315,80,0,1),\\
u={}&(171,133,19,190,266,95,76,76,56,0,1),\\
w={}&(342,247,342,171,266,38,190,75,213,0,281,1).
\end{align*}
Their exact modular identities are
\begin{equation}\label{eq:kernel}
 vA_{14}\equiv0\pmod{729},\qquad
 uA_{10}\equiv0\pmod{361},\qquad
 wA_{11}\equiv0\pmod{361}.
\end{equation}
Each identity means simply that, for every column $k$,
$\sum_{n=k}^d v_n\binom nk^2\binom{n+k}{k}^2$ is divisible by the stated
modulus, and similarly for $u,w$. There are only 38 column checks.
The verification code below specifies them without a matrix inversion.

For a vector $z=(z_0,\ldots,z_d)$, prime $p$, and integer $a\ge0$, define
\[
 U_{z,p}(a)=\sum_{n=0}^d z_n
       \sum_{\substack{0\le k\le n\\p\nmid B_{n,k}}}B_{n,k}^{\,a}.
\]
If $r\ge6$, all terms with $3\mid B_{n,k}$ vanish modulo $729=3^6$.
Terms with $19\mid B_{n,k}$ vanish modulo $361=19^2$ already for $r\ge2$.
Thus~\eqref{eq:kernel} forces, for $r\ge6$,
\begin{equation}\label{eq:units}
 U_{v,3}(r)\equiv0\pmod{729},\qquad
 U_{u,19}(r)\equiv U_{w,19}(r)\equiv0\pmod{361}.
\end{equation}

\subsection*{The condition at 3}
For $3\nmid b$, Euler's theorem gives $b^{18}\equiv1\pmod{27}$.
Define $q_3(b)=(b^{18}-1)/27$ modulo 27. Direct integer arithmetic gives
\begin{equation}\label{eq:derivative}
 \sum_{n=0}^{14}v_n
 \sum_{\substack{0\le k\le n\\3\nmid B_{n,k}}}
 B_{n,k}^{\,h}q_3(B_{n,k})\equiv0\pmod{27}
 \quad(0\le h<18).
\end{equation}
As $27^2=729$, the binomial theorem now proves, for every $t\ge0$,
\[
 U_{v,3}(h+18t)\equiv U_{v,3}(h)\pmod{729}.
\]
The residues $U_{v,3}(h)$ are all divisible by 81. Their quotients modulo 9
are the following complete table:
\[
\begin{array}{c|rrrrrrrrr}
h&0&1&2&3&4&5&6&7&8\\\hline
U_{v,3}(h)/81&4&2&6&4&2&0&1&8&3\\[2pt]
h&9&10&11&12&13&14&15&16&17\\\hline
U_{v,3}(h)/81&1&8&6&7&5&0&7&5&3
\end{array}
\]
Here and below division of a residue by 81 or 19 means first taking its
representative in $[0,m)$, which is divisible by that integer.
The first necessary congruence in~\eqref{eq:units} therefore forces
\begin{equation}\label{eq:r18}
 r\equiv5\ \hbox{or}\ 14\pmod{18}.
\end{equation}

\subsection*{Two incompatible conditions at 19}
Write $r=h+18t$ with $h\in\{5,14\}$ and $t\ge0$ as permitted
by~\eqref{eq:r18}. For $19\nmid b$, define
$q_{19}(b)=(b^{18}-1)/19$ modulo 19. Fermat's theorem and the binomial
theorem give
\[
 b^{h+18t}\equiv b^h(1+19tq_{19}(b))\pmod{361}.
\]
For $z=u,w$ and these two $h$, $U_{z,19}(h)$ is divisible by 19, so
\[
 \frac{U_{z,19}(h+18t)}{19}\equiv a_{z,h}+b_{z,h}t\pmod{19},
\]
where $a_{z,h}=U_{z,19}(h)/19$ modulo 19 and
\[
 b_{z,h}=\sum_n z_n\sum_{\substack{k\le n\\19\nmid B_{n,k}}}
 B_{n,k}^{\,h}q_{19}(B_{n,k})\pmod{19}.
\]
The complete table is
\[
\begin{array}{cc|rr|r}
z&h&a_{z,h}&b_{z,h}&\text{required }t\pmod{19}\\\hline
u&5 &3&2&8\\
w&5 &7&14&9\\
u&14&9&10&1\\
w&14&1&5&15
\end{array}
\]
The two required residues of $t$ disagree for either value of $h$.
This contradicts~\eqref{eq:units} and excludes every $r\ge6$.

\subsection*{The remaining exponents}
Forward substitution in~\eqref{eq:def}, with $s=2$, gives
\[
 c_3^{(3,2)}=\frac{343}{5},\qquad
 c_3^{(4,2)}=\frac{11612}{5},\qquad
 c_3^{(5,2)}=\frac{342076}{5}.
\]
None is an integer. Together with the preceding contradiction, this
proves Theorem~\ref{thm:main}.\qed

\subsection*{Exact finite verification}
All displayed data use only binomial coefficients, modular exponentiation,
and finite sums. The accompanying \texttt{verify\_certificate.py} checks
the 38 kernel columns, all 18 identities~\eqref{eq:derivative}, both tables,
and the three exceptional fractions. To make the certificate computation
fully specified in the paper as well, its arithmetic core is:
\begin{verbatim}
from math import comb
def B(n,k): return comb(n,k)*comb(n+k,k)
def terms(z,p):
    return [(a,B(n,k)) for n,a in enumerate(z)
            for k in range(n+1) if B(n,k)%p]
def U(z,p,h,m):
    return sum(a*pow(b,h,m) for a,b in terms(z,p))%m
def kernel(z,m):
    return [sum(z[n]*B(n,k)**2 for n in range(k,len(z)))%m
            for k in range(len(z))]
def slope(z,p,h,divisor,m):
    return sum(a*pow(b,h,m//divisor)*
               ((pow(b,18,m)-1)//divisor)
               for a,b in terms(z,p))%(m//divisor)
# Insert the three vectors v,u,w displayed above.
# kernel(v,729), kernel(u,361), kernel(w,361) are zero lists.
# For h=0,...,17, slope(v,3,h,27,729) is zero;
# U(v,3,h,729)//81 gives the first table.
# For z=u,w and h=5,14, the second table is
# (U(z,19,h,361)//19, slope(z,19,h,19,361)).
\end{verbatim}
The standalone script additionally checks the complete Euler periods
486 and 342 by direct enumeration and reproduces the source's six
small examples. These supplementary checks agree with the shorter
algebraic certificates above. The proof covers all exponents by the
binomial identities; it does not infer an infinite statement from a
bounded search. No independent peer review or formal proof-assistant
verification is claimed.

\section*{Four further research tasks}
\paragraph{Task 2.} Determine the least $N$ for which no $r>2$ makes
$c_0^{(r,2)},\ldots,c_N^{(r,2)}$ all integers. The theorem gives $N\le14$;
an exact admissible-exponent classification for $N\le13$ would settle
minimality.
\paragraph{Task 3.} For each fixed $s\ge3$, determine whether a finite
$N(s)$ excludes every $r>s$, and obtain explicit certificates whenever
such an $N(s)$ exists.
\paragraph{Task 4.} For fixed $s,N$, construct the complete eventually
periodic set of exponents giving integral initial coefficients, with a
proved onset and period derived from the rational triangular inverse.
\paragraph{Task 5.} Decide Conjecture 4.2 of~\cite{GZ} for all $r>s>1$:
must some coefficient be nonintegral? The present theorem resolves its
$s=2$ specialization, but does not resolve all $s>1$.

\begin{thebibliography}{9}
\bibitem{GZ} V. J. W. Guo and J. Zeng,
\emph{On Zudilin's $q$-question about Schmidt's problem},
Electronic Journal of Combinatorics \textbf{19}(3) (2012), P4,
\href{https://doi.org/10.37236/2231}{doi:10.37236/2231}.
Conjectures 4.1 and 4.2 appear on page 4 of the final article.
The current author manuscript is
\href{https://arxiv.org/abs/1204.0187v3}{arXiv:1204.0187v3}.
\end{thebibliography}
\end{document}
